%!TEX root = ../main.tex

\chapter{Background}\label{ch:background}

This chapter reviews the technical foundations on which the experiments in this
thesis are built. Section~\ref{sec:bg:llm} introduces \glspl{llm}, their strengths,
and their limitations as standalone question-answering systems. Section~\ref{sec:bg:rag}
describes \gls{rag} as a mechanism for grounding generation in external documents.
Section~\ref{sec:bg:cqa} discusses the specific challenges of conversational
settings. Section~\ref{sec:bg:kg} briefly introduces knowledge graphs as a
structured complement to document-based retrieval.


% ===========================================================================
\section{Large Language Models}\label{sec:bg:llm}

\Glspl{llm} are neural language models with large numbers of parameters, typically
trained on broad web-scale corpora using a self-supervised next-token prediction
objective~\cite{brown2020gpt3}. The dominant architecture is the
transformer~\cite{vaswani2017attention}, which uses multi-head self-attention to
model dependencies between all positions in a sequence regardless of distance.

After pre-training, \glspl{llm} are commonly adapted through \emph{instruction
tuning}: fine-tuning on a curated collection of (instruction, response) pairs,
optionally followed by reinforcement learning from human feedback. This teaches the model to follow instructions, produce well-formatted
answers and, to some extent, express uncertainty.

\subsection{In-context Learning and Few-Shot Prompting}

An important property of large \glspl{llm} is \emph{in-context learning}: the
model's behaviour can be steered at inference time by providing a few
(input, output) examples in the prompt, without any gradient updates. This
\emph{few-shot prompting} approach is used in this work to make the model answer in
plain text, in the second person and in one or two sentences, and to show it when to
abstain.

\subsection{Limitations}

Despite their fluency, \glspl{llm} exhibit several well-documented
limitations relevant to document-grounded \gls{qa}:

\begin{description}
    \item[Hallucination.] Models frequently generate plausible-sounding but
        factually incorrect statements, particularly when queried about specific
        documents, dates, or numerical values not seen frequently during
        training~\cite{lewis2020rag}.

    \item[Static knowledge.] Parametric knowledge is frozen at training time.
        Domain-specific or recently updated documents cannot be incorporated
        without retraining or fine-tuning.

    \item[No grounding signal.] Without explicit access to source documents,
        the model cannot cite evidence, and the user has no mechanism to verify
        the origin of a claim.

    \item[Unanswerable question handling.] Without retrieval, models tend to
        answer even questions that lack supporting evidence, introducing
        uncontrolled hallucination on unanswerable inputs.
\end{description}

These limitations motivate the \gls{rag} architecture described in the next section.


% ===========================================================================
\section{Retrieval-Augmented Generation}\label{sec:bg:rag}

\Gls{rag}, proposed by Lewis~et~al.~\cite{lewis2020rag}, addresses the limitations
of purely parametric models by augmenting each query with passages retrieved from
an external document store. The retriever selects $k$ passages most relevant to the
query, and the generator produces an answer conditioned on both the query and the
retrieved passages. This decoupling allows the knowledge base to be updated without
retraining the language model.

\subsection{Sparse Retrieval: BM25}

Before the advent of neural retrieval, the dominant retrieval model for information
retrieval was based on term frequency and inverse document frequency.
\gls{bm25}~\cite{robertson2009bm25} is a probabilistic extension of TF-IDF that
introduces document length normalisation and a saturation parameter:
\begin{equation}
    \mathrm{BM25}(q, d) = \sum_{t \in q}
        \mathrm{IDF}(t) \cdot
        \frac{f(t,d) \cdot (k_1 + 1)}
             {f(t,d) + k_1 \bigl(1 - b + b \cdot |d| / \overline{|d|}\bigr)},
\end{equation}
where $f(t, d)$ is the term frequency of term $t$ in document $d$, $|d|$ is the
document length, $\overline{|d|}$ is the average document length in the corpus,
and $k_1$, $b$ are hyperparameters.
BM25 is fast, interpretable, and effective for keyword-rich queries, but it cannot
match documents on the basis of semantic similarity and is sensitive to vocabulary
mismatch.

\subsection{Dense Retrieval}

Dense retrieval methods~\cite{karpukhin2020dpr} encode queries and documents as
dense vectors using a neural encoder (typically a bi-encoder based on a pre-trained
transformer), and retrieve by approximate nearest-neighbour search in the embedding
space. This allows the retriever to handle paraphrase and semantic variation that
would cause keyword-based methods to fail.

This work uses \texttt{BAAI/bge-small-en-v1.5}~\cite{bge2023}, a 384-dimensional
bi-encoder fine-tuned for retrieval tasks. Section-level embeddings are pre-computed
once and stored on disk; query embeddings are computed at inference time and matched
by cosine similarity.

\subsection{Hybrid Retrieval and Reciprocal Rank Fusion}

Sparse and dense retrievers have complementary strengths: BM25 excels at exact
keyword matching while dense retrievers handle semantic variation. Hybrid retrieval
combines their ranked lists using \gls{rrf}~\cite{cormack2009rrf}:
\begin{equation}
    \mathrm{RRF}(d) = \sum_{r \in \mathcal{R}}
        \frac{1}{k_{\mathrm{rrf}} + \mathrm{rank}_r(d)},
\end{equation}
where $\mathcal{R}$ is the set of ranked lists and $k_{\mathrm{rrf}}$ is a constant
(set to $60$ in the experiments). The fused score requires no normalisation and
tends to rank documents that appear highly in multiple lists above those that rank
highly in only one.

\subsection{Cross-Encoder Reranking}

A limitation of the bi-encoder architecture is that query and document are encoded
independently, without cross-attention between them. \emph{Cross-encoder}
rerankers~\cite{nogueira2020passage} concatenate the query and document as a single
input to a transformer, enabling full cross-attention. This produces more accurate
relevance scores at the cost of higher computational complexity.

This work uses \texttt{BAAI/bge-reranker-base} as a second stage:
the retriever fetches the top~50 candidates, the reranker scores all 50
(query, section) pairs, and the top~5 are passed to the generator.
Section texts exceeding 512~tokens (the model's maximum) are silently truncated by
the model.


% ===========================================================================
\section{Conversational Question Answering}\label{sec:bg:cqa}

\Gls{cqa} extends single-turn \gls{qa} to multi-turn dialogues, where the meaning
of each question may depend on the dialogue history. This creates two challenges
not present in the single-turn setting.

\subsection{Query Decontextualisation}

Many follow-up questions in a dialogue are \emph{decontextualised}: they use
pronouns, ellipsis, or topic continuations that are only interpretable in the
context of prior turns. For example, the question ``What are the fees?'' following
a discussion of a driver's licence renewal requires a reader (or retrieval system)
to resolve ``fees'' as ``driver's licence renewal fees'' to formulate an effective
retrieval query.

Query rewriting addresses this by generating a self-contained reformulation of the
follow-up question given the dialogue history~\cite{yu2020few}. Here it is implemented
as an additional \gls{llm} call that produces a rewritten query before retrieval,
using a sliding window of the last six dialogue turns as context.

\subsection{Document Switch}

In multi-document dialogues, successive turns may concern different source
documents. A \emph{document-switch} turn requires the retrieval system to locate
the new relevant document without being misled by the strong lexical signal of the
previous document's context. Query rewriting can both help and hurt here: it
improves same-document follow-ups but can encode the previous document's vocabulary
so strongly that retrieval is pulled towards the old document rather than the new one.

\subsection{Unanswerable Questions}

In realistic deployments, users will ask questions for which the corpus contains
no answer. A reliable system must \emph{abstain} on such questions rather than
generating a fluent but unsupported answer. This creates a calibration problem:
abstention should fire on genuinely unanswerable questions but not on answerable
ones (false abstention). This is measured by abstain recall, false-abstention rate
and abstain precision.


% ===========================================================================
\section{Knowledge Graphs}\label{sec:bg:kg}

\Glspl{kg} represent information as a set of typed triples
$(e_1, r, e_2)$, where $e_1$ and $e_2$ are entities and $r$ is a named relation.
This structured representation captures explicit relational knowledge that is
implicit in unstructured text, enabling precise graph-query retrieval and
multi-hop reasoning.

The integration of \glspl{kg} with \glspl{llm}~\cite{pan2023unifying} is an active
research direction. In the context of \gls{rag}, a \gls{kg} constructed from the
document corpus can serve as an additional retrieval source: graph traversal from
entities mentioned in the query can surface supporting facts that dense retrieval
over raw text would miss, particularly for multi-hop or relation-centric questions.

Chapter~\ref{ch:kg} describes the construction of a \gls{kg} from the \gls{mdd}
corpus, following the methodology of~\textcite{negro2026kg}, and
Chapter~\ref{ch:kgret} its use for retrieval.
